\documentclass[a4paper,11pt]{article}
\usepackage[brazilian]{babel}
\usepackage[utf8]{inputenc}
\usepackage[T1]{fontenc}
\usepackage{epsfig}
\usepackage{array}

\usepackage{graphicx, url}
\usepackage[usenames]{color}
\usepackage{rotating}
\usepackage{url}
\usepackage{subfigure}
\usepackage{xspace}
\usepackage[verbose, dvips=true, pdftex=false, vtex=false, paperwidth=210mm,paperheight=295mm,top=35mm,left=30mm,bottom=35mm,right=30mm]{geometry}
\usepackage[authoryear]{natbib}
\usepackage{fancyhdr}
\usepackage{setspace}


\DeclareGraphicsExtensions{.png,.jpg,.pdf,.mps}
\bibliographystyle{icmc2}

\clubpenalty=10000
\widowpenalty=10000
\newcommand{\rev}[1]{{\color{red}#1}}

\newcommand{\titlehd}{\sffamily\bfseries\selectfont \scriptsize Uma Abordagem de Reestruturação de Sistemas Baseada em Requisitos de Qualidade Pré-Estabelecidos }

\pagestyle{fancy}
\headheight 61pt
\fancyhead{}
\fancyhead[LO]{\epsfig{file=icmc.jpg,height=15mm}}
\fancyhead[RO]{}
\fancyfoot{}

\begin{document}

\begin{center}
\vspace*{1cm}
{\fontsize{17.28}{20}\sffamily\bfseries\selectfont Uma Abordagem de Reestruturação de Sistemas Baseada em Requisitos de Qualidade Pré-Estabelecidos \\\ \\

\small{RELATÓRIO CIENTÍFICO PARCIAL - 01/06/2012 a 15/03/2013}\\ \
\textsf{Apresentado à Fundação de Amparo à Pesquisa do Estado de São Paulo - FAPESP}
 \\
\vspace{3cm}}


\begin{table}[!th]
  \begin{center}
    \begin{tabular}{ll}

	  Número do Processo:      &   FAPESP 2012/05168-4     \\
      Período:        &  Junho/2012 a Março/2013                    \\
	                   &                                \\
      Bolsista:           &  Rafael Serapilha Durelli (rdurelli@icmc.usp.br)                                     \\
	  
	  \vspace{0.6cm}
      Orientador:        &  Prof. Dr. Márcio E. Delamaro (delamaro@icmc.usp.br)                                           \\

	         &                              \\
	          &  \\
                       &                                  \\
       &                                       \\
              &                                             \\ \\
      \vspace{1cm}

    \end{tabular}
  \end{center}
\end{table}

\vfill
{\fontsize{18}{20}\sffamily\bfseries\selectfont USP - São Carlos \\ Janeiro de 2012}

\end{center}

\vspace{3cm}

\pagebreak
\fancyhead[RO]{\titlehd}
\fancyfoot{\centerline{\thepage}}
\setcounter{page}{1}

\parskip 6pt

\doublespacing
\begin{abstract}
 Relatório Científico Parcial apresentado à FAPESP (Fundação de Amparo à Pesquisa do Estado de São Paulo) com o objetivo de elucidar as atividades realizadas  
 pelo bolsista Rafael Serapilha Durelli durante o primeiro período de vigência da bolsa concebida sob o Processo Número 2012/05168-4. 
 O referido período teve inicio em Junho de 2012 e foi finalizado em Março de 2013. 
 Além disso, este relatório também descreve as atividades que estão em andamentos, bem como as atividades a serem realizadas no próximo período. 
 Vale ressaltar que a partir de resultados preliminares alguns artigos científicos foram elaborados e publicados.
%Software é um produto que tende a evoluir constantemente com o intuito de satisfazer às necessidades de seus usuários. Para isso é necessário submetê-lo a constantes atividades de manutenção, que podem ocasionar um possível degradação no código-fonte. Dessa forma, o objetivo desta proposta de doutorado é o desenvolvimento de uma abordagem que auxilie o engenheiro de software a escolher o melhor conjunto de transformações para um sistema legado com base em requisitos de qualidade pré-estabelecidos e com isso possa identificar o impacto que tais transformações irão ocasionar no sistema legado. Isto é, a abordagem deverá indicar conjuntos de possíveis transformações para a evolução do sistema legado, além disso, deverá informar ao engenheiro o quanto tais transformações irão afetar a evolução/estrutura desse sistema. Com isso, o engenheiro de software poderia previamente escolher qual o melhor conjunto de  transformação se aplica dado um requisito de qualidade pré-estabelecido. Após a escolha das transformações a abordagem também deverá realizar automaticamente tais transformações a fim de reestruturar o sistema legado em um novo sistema. Será avaliado a abordagem por meio de experimentos controlados.

\end{abstract}


      



\section{Introdução}
	 Este relatório tem por objetivo apresentar as atividades realizadas pelo bolsista Rafael Serapilha Durelli durante o período de Junho/2012 a Março/2013,
 referente à bolsa de doutorado concebida pela Fundação de Amparo à Pesquisa do Estado de São Paulo (FAPESP) sob o Processo Número 2012/05168-4. 
 
 É importante salientar que o trabalho em questão tem sido desenvolvido no Departamento de Ciências da Computação e Estatística do Instituto de Ciência 
 Matemáticas e de Computação (ICMC) da Universidade de São Paulo (campos São Carlos/SP).  
 Este trabalho se insere no contexto do grupo de pesquisas em Engenheira de Software, sob a orientação do Prof. Dr. Márcio Eduardo Delamaro.  
 Além disso, este trabalho esta sendo executado em colaboração com o grupo de engenharia de software da Universidade Federal de São Carlos (UFSCAR)\footnote{http://dc.ufscar.br}.  
 Mais especificamente em colaboração com o Prof. Dr. Valter Vieira de Camargo\footnote{http://buscatextual.cnpq.br/buscatextual/visualizacv.do?id=S819089},
 o qual tem grande experiência na área de engenharia de software com ênfase no desenvolvimento de frameworks no contexto da programação orientada a aspectos e reuso de software. 

Por fim, ressalta-se que o bolsista criou um vínculo científico com o \textit{Institut National de Recherche en Informatique et en Automatique} (INRIA),  
onde irá realizar um ano de doutorado sanduíche sobre orientação do  
 Prof. Dr. Nicolas Anquetil\footnote{http://rmod.lille.inria.fr/web/pier/team/Nicolas-Anquetil} o qual tem grande experiência na área de manutenção e reengenharia de software.

\section{Plano de Trabalho}\label{plano_trabalho}

	\subsection{Resumo do Plano Inicial}\label{resumo_plano_inicial}

	\subsection{Cronograma}\label{cronograma}

\section{Resumo das Atividades Realizadas no Período}\label{atividades_realizadas}

	Nesta seção são apresentadas e detalhadas as atividades realizadas durante o primeiro período de bolsa,  
 correspondente aos meses de Junho/2012 a Março/2013.
 Em resumo, foram cursadas um conjunto de disciplinas para a obtenção dos créditos necessários do Programa de Pós-Graduação em Ciências de Computação e Matemática Computacional do ICMC. Além disso, estudou-se a abordagem Modernização Dirigida à Arquitetura
  ( em inglês - \textit{Architecture-Driven Modernization} - ADM), bem como seu principal ativo, o Meta-Modelo de Descoberta de Conhecimento (em inglês - \textit{Knowledge Discovery Meta-Model} - KDM). 

Realizou-se também uma revisão bibliográfica sistemática para investigar as principais técnicas e trabalhos disponíveis na literatura para auxiliar o desenvolvido do projeto em questão. O bolsista elaborou a monografia para o Exame de Qualificação e apresentou a mesma para a banca. Também são descritas nesta seção as publicações e eventos científicos nos quais o bolsista participou diretamente.

	\subsection{Atividades Técnicas Previstas no Cronograma do Projeto}

		\subsection{Disciplinas Regulares do Curso de Doutorado do ICMC/USP} % (fold)
\label{sub:disciplinas_regulares_do_curso_de_doutorado_do_icmc_usp}

De acordo com o tema central do projeto, foram escolhidas algumas disciplinas julgadas importantes para o bom aproveitamento do trabalho proposto no cronograma inicial. O quadro de disciplina cursadas pelo bolsista é mostrado a seguir:
% subsection disciplinas_regulares_do_curso_de_doutorado_do_icmc_usp (end)

\begin{enumerate}
	\item Engenharia de Software Experimental
		\begin{itemize}
			\item Carga Horária - 90
			\item Créditos: 6 (seis)
			\item Conceito Obtido: A
		\end{itemize}
	\item Revisão Sistemática em Engenharia de Software
		\begin{itemize}
			\item Carga Horária - 90
			\item Créditos: 6 (seis)
			\item Conceito Obtido: A
		\end{itemize}
	\item Validação e Teste de Software
		\begin{itemize}
			\item Carga Horária - 180
			\item Créditos: 12 (doze)
			\item Conceito Obtido: A
		\end{itemize}
	\item Preparação Pedagógica
		\begin{itemize}
			\item Carga Horária - 60
			\item Créditos: 4 (quatro)
			\item Conceito Obtido: A 
		\end{itemize}
	\item Especificação formal de software
		\begin{itemize}
			\item Carga Horária - 180
			\item Créditos: 12 (doze)
			\item Conceito Obtido: A
		\end{itemize}
\end{enumerate}

De acordo com as regras do Programa de Pós-Graduação em Ciências de Computação e Matemática Computacional do ICMC os alunos devem obter uma quantidade mínima de créditos para realizar determinadas tarefas, tais como, realizar o exame de qualificação e entregar a dissertação/tese. Nesse contexto, vale ressaltar que o bolsista em questão com o comprimento das disciplinas listadas a cima obteu-se um total de 40 créditos, sendo apto para realizar o exame de qualificação que exigi o mínimo de 24 créditos.%, bem como 36 créditos são necessários para o depósito da tese. 

Maiores informações sobre o desempenho do bolsista em cada uma das disciplinas e os conceitos obtidos podem ser avaliados mais cuidadosamente no histórico escolar deste relatório (ver Apêndice A).

%Conforme mencionado, durante o período deste relatório, o bolsista realizou tanto atividades técnicas requeridas para a concretização do seu projeto de doutorado quanto atividades exigidas pelo Programa de Pós-Graduação. Neste contexto, nesta subseção são apresentas e detalhadas as atividades técnicas realizadas até o momento. A seguir, cada uma das atividades realizadas é descrita.

\subsection{Estudo sobre Reengenharia de Software} % (fold)
\label{sub:reengenharia}


Software é um produto que evolui constantemente para satisfazer às necessidades de seus usuários. Para isso é necessário submetê-lo a constantes atividades de manutenção, que podem degradar o código-fonte, tornando-se cada vez mais difícil de mantê-lo uma vez que, na maioria das vezes, o software não é atulizado, bem como a sua a documentação, culminando em uma situação em que a única documentação confiável é o código-fonte. Sistemas com essas características são denominados sistemas legados. Usualmente, sistemas legados são candidatos à reengenharia.


Segundo~\citet{refactImpro} reengenharia no contexto da evolução do software é utilizada para melhorar a qualidade do software, ou seja, melhorar a extensibilidade, modularidade, reusabilidade, complexidade e manutenibilidade. Por exemplo, reengenharia é útil para converter sistemas legados ou códigos deteriorados em unidades mais modularizadas ou até mesmo migrar tais sistemas para diferentes linguagens de programação e paradigmas.


De acordo com~\citet{chikofskyTax}, a reengenharia de software  é necessário para converter o código legado ou deteriorado em um código mais modular e estruturado, ou até mesmo alterar o paradigma de programação. Segundo~\citet{Sneed:2005} mais da metade dos projetos que aplicam reengenharia falham ao lidar com desafios específicos. De acordo com esse autor, tanto a carência de  padronização durante a atividade de reengenharia quanto a falta de apoio computacional efetivo são os principais problemas que acarretam esse grande número de falhas.  

A falta de um processo padronizado de reengenharia usualmente é um problema pois, usualmente tal processo é realizado de forma totalmente \textit{ad hoc}. Além disso, a carência de um apoio computacional efetivo para auxiliar a atividade de reengenharia também é um problema, uma vez que tal atividade não é trivial e requer que várias mudanças sejam realizadas tanto no código-fonte como em outros artefatos, e.g., documentos de requisitos, casos de uso e diagrama de classes. 

% subsection evolu_o_de_software_e_reengenharia (end)

\subsection{Estudo sobre Desenvolvimento Dirigido a Modelos} % (fold)
\label{sub:desenvolvimento_dirigido_a_modelos}

A proposta do Desenvolvimento Dirigido a Modelos (\textit{Model-Driven Development - MDD}) é reduzir a distância semântica entre o problema do domínio e a solução/implementação. Assim, o engenheiro de software não precisa interagir inteiramente com o código-fonte, podendo-se concentrar em modelos que possuem maiores níveis de abstração. Um mecanismo é responsável por gerar automaticamente o código-fonte por meio dos modelos. No MDD, modelos não apenas guiam as tarefas de desenvolvimento e manutenção, mas são partes integrante do software assim como o código-fonte, servindo como entrada para ferramentas de geração de código reduzindo os esforços dos desenvolvedores~\cite{Bittar, Kleppe:2003}


No desenvolvimento tradicional, ou seja, sem seguir o MDD, artefatos de alto nível (ex., modelos, diagramas) usualmente são produzidos antes da codificação e costumam ser úteis apenas nas etapas iniciais do ciclo de desenvolvimento. Conforme o desenvolvimento evolui, mudanças são aplicadas somente no código-fonte e não nos modelos/diagramas. Assim, tais artefatos acabam se tornando incoerentes, ou seja, não refletem o que o código-fonte apresenta, o que faz com que o tempo e os esforços gastos na construção desses artefatos não sejam diretamente aproveitados na produção do software~\cite{Bittar}. Contrapartida, o MDD tem o foco nos modelos e busca simplificar o processo de desenvolvimento de software. Modelos são mais intuitivos para representação do conhecimento e menos dependentes do código-fonte, de forma que podem ser reutilizados facilmente em diferentes projetos. Diferentemente, o código-fonte possui uma linguagem que é densa e codificada, tornando-se difícil identificar, extrair e reutilizar o conhecimento apenas pela leitura do mesmo~\citep{Kleppe:2003}.



% subsection desenvolvimento_dirigido_a_modelos (end)

\subsection{Estudo sobre Modernização Dirigida à Arquitetura } % (fold)
\label{sub:estudo_sobre_moderniza_o_dirigida_arquitetura_}

Reengenharia de Software e MDD são duas abordagens diferentes e que foram pesquisadas separadamente durante anos. Porém, recentemente pesquisadores identificam interesses análogos entre ambas abordagens. 

Como dito anteriormente, o principal objetivo da Reengenharia de Software é converter sistemas legados para novos sistemas sem alterar a sua funcionalidade, ou seja, isso implica em entender como o sistema legado foi implementado. Portanto, tal sistema legado deve ser convertido a um nível maior de abstração com o intuito de auxiliar o engenheiro de software a se concentrar apenas em informações importantes desse sistema, nesse ponto que MDD e Reengenharia de Software podem ser utilizados em conjunto.

Nesse contexto, a OMG OMG (\textit{Object Management Group}) propôs a abordagem Modernização Dirigida à Arquitetura (\textit{Architecture-Driven Modernization} - ADM), a qual tem como intuito automatizar e formalizar/padronizar o problema tradicional da Reengenharia de Software, i.e., ADM melhora a abordagem  de reengenharia de software tradicional com a utilização de MDD~\cite{PerezCastillo:2011jo}. Segundo~\citet{rezCastillo:2011gm} ADM é uma padronização definida pela OMG para auxiliar a atividade de reengenharia de sistemas legados, com o diferencial da utilização dos princípios da MDD (ver Seção~\ref{Cap2_Sec2_Desenvolvimento_Dirigido_a_Modelos}). ADM difere das abordagens tradicionais de reengenharia de software por dois principais motivos: (\textit{i}) ADM considera todos os artefatos de um sistema legado como modelos e (\textit{ii}) refatorações do sistema legado são realizadas nos modelos e depois gera-se um novo código-fonte refatorado tendo como base tais modelos. 

De acordo com~\citet{PerezCastillo:2011jo} ADM tem algumas vantagens quando comparada com abordagens de reengenharia convencionais: (\textit{i}) permite que pesquisadores definam técnicas de refatorações independente de linguagem e plataforma; (\textit{ii}) refatorações podem ser definidas como modelos, assim, podem ser reutilizadas, ou seja, utiliza o metamodelo KDM, o qual tem como um dos objetivos aumentar a interoperabilidade entre as ferramentas de reengenharia e (\textit{iii}) com a abordagem ADM é possível criar técnicas de refatorações genéricas e específicas. 

% subsection estudo_detalhado_sobre_moderniza_o_dirigida_arquitetura_ (end)

\begin{itemize}

\item estudar ferramentas para criar uma Linguagem Específica de Domínio: fez-se necessário identificar e estudar algumas ferramentas que auxiliam o desenvolvimento de Linguagem Específica de Domínio. Foram identificadas as seguintes ferramentas, XText\footnote{http://www.eclipse.org/Xtext/}, Eclipse EMF\footnote{http://www.eclipse.org/modeling/emf/} e Eclipse GMF\footnote{http://www.eclipse.org/modeling/gmp/}. Tais ferramentas serão estudadas para o desenvolvimento de um Linguagem Específica de Domínio, a qual será utilizada no terceiro passo da abordagem proposta;

\item Estudo detalhado da abordagem denominada Arquitetura Dirigida a Modelo - (ADM): conforme descrito no projeto submetido anteriormente pretende-se criar uma abordagem para auxiliar o engenheiro de software durante a atividade de reestruturação de sistemas legados. Dessa forma, vez se necessário estudar realizar a modernização de sistemas legados com a utilizanção de 

\end{itemize}


	\subsection{Disciplinas Regulares do Curso de Doutorado do ICMC/USP} % (fold)
		\label{sub:disciplinas_regulares_do_curso_de_doutorado_do_icmc_usp}
		 
		 
De acordo com o tema central do projeto, foram escolhidas algumas disciplinas julgadas importantes para o bom aproveitamento do trabalho proposto no cronograma inicial. O quadro de disciplina cursadas pelo bolsista é mostrado a seguir:
% subsection disciplinas_regulares_do_curso_de_doutorado_do_icmc_usp (end)

\begin{enumerate}
	\item Engenharia de Software Experimental
		\begin{itemize}
			\item Carga Horária - 90
			\item Créditos: 6 (seis)
			\item Conceito Obtido: A
		\end{itemize}
	\item Revisão Sistemática em Engenharia de Software
		\begin{itemize}
			\item Carga Horária - 90
			\item Créditos: 6 (seis)
			\item Conceito Obtido: A
		\end{itemize}
	\item Validação e Teste de Software
		\begin{itemize}
			\item Carga Horária - 180
			\item Créditos: 12 (doze)
			\item Conceito Obtido: A
		\end{itemize}
	\item Preparação Pedagógica
		\begin{itemize}
			\item Carga Horária - 60
			\item Créditos: 4 (quatro)
			\item Conceito Obtido: A 
		\end{itemize}
	\item Especificação formal de software
		\begin{itemize}
			\item Carga Horária - 180
			\item Créditos: 12 (doze)
			\item Conceito Obtido: A
		\end{itemize}
\end{enumerate}

De acordo com as regras do Programa de Pós-Graduação em Ciências de Computação e Matemática Computacional do ICMC os alunos devem obter uma quantidade mínima de créditos para realizar determinadas tarefas, tais como, realizar o exame de qualificação e entregar a dissertação/tese. Nesse contexto, vale ressaltar que o bolsista em questão com o comprimento das disciplinas listadas acima obteu-se um total de 40 créditos, sendo apto para realizar o exame de qualificação que exigi o mínimo de 24 créditos.%, bem como 36 créditos são necessários para o depósito da tese. 

Maiores informações sobre o desempenho do bolsista em cada uma das disciplinas e os conceitos obtidos podem ser avaliados mais cuidadosamente no histórico escolar deste relatório (ver Apêndice A).

%Conforme mencionado, durante o período deste relatório, o bolsista realizou tanto atividades técnicas requeridas para a concretização do seu projeto de doutorado quanto atividades exigidas pelo Programa de Pós-Graduação. Neste contexto, nesta subseção são apresentas e detalhadas as atividades técnicas realizadas até o momento. A seguir, cada uma das atividades realizadas é descrita.
	
	\subsection{Estudo sobre Reengenharia de Software} % (fold)
	\label{sub:estudo_sobre_reengenharia_de_software}
		 
		 
Software é um produto que evolui constantemente para satisfazer às necessidades de seus usuários. Para isso é necessário submetê-lo a constantes atividades de manutenção, que podem degradar o código-fonte, tornando-se cada vez mais difícil de mantê-lo uma vez que, na maioria das vezes, o software não é atulizado, bem como a sua a documentação, culminando em uma situação em que a única documentação confiável é o código-fonte. Sistemas com essas características são denominados sistemas legados. Usualmente, sistemas legados são candidatos à reengenharia.


Segundo~\citet{refactImpro} reengenharia no contexto da evolução do software é utilizada para melhorar a qualidade do software, ou seja, melhorar a extensibilidade, modularidade, reusabilidade, complexidade e manutenibilidade. Por exemplo, reengenharia é útil para converter sistemas legados ou códigos deteriorados em unidades mais modularizadas ou até mesmo migrar tais sistemas para diferentes linguagens de programação e paradigmas.


De acordo com~\citet{chikofskyTax}, a reengenharia de software  é necessário para converter o código legado ou deteriorado em um código mais modular e estruturado, ou até mesmo alterar o paradigma de programação. Segundo~\citet{Sneed:2005} mais da metade dos projetos que aplicam reengenharia falham ao lidar com desafios específicos. De acordo com esse autor, tanto a carência de  padronização durante a atividade de reengenharia quanto a falta de apoio computacional efetivo são os principais problemas que acarretam esse grande número de falhas.  

A falta de um processo padronizado de reengenharia usualmente é um problema pois, usualmente tal processo é realizado de forma totalmente \textit{ad hoc}. Além disso, a carência de um apoio computacional efetivo para auxiliar a atividade de reengenharia também é um problema, uma vez que tal atividade não é trivial e requer que várias mudanças sejam realizadas tanto no código-fonte como em outros artefatos, e.g., documentos de requisitos, casos de uso e diagrama de classes. 

% subsection evolu_o_de_software_e_reengenharia (end)
		
	 \subsection{Estudo sobre Desenvolvimento Dirigido a Modelos} % (fold)
	 	\label{sub:estudo_sobre_desenvolvimento_dirigido_a_modelos}
	 	
	 	

A proposta do Desenvolvimento Dirigido a Modelos (\textit{Model-Driven Development - MDD}) é reduzir a distância semântica entre o problema do domínio e a solução/implementação. Assim, o engenheiro de software não precisa interagir inteiramente com o código-fonte, podendo-se concentrar em modelos que possuem maiores níveis de abstração. Um mecanismo é responsável por gerar automaticamente o código-fonte por meio dos modelos. No MDD, modelos não apenas guiam as tarefas de desenvolvimento e manutenção, mas são partes integrante do software assim como o código-fonte, servindo como entrada para ferramentas de geração de código reduzindo os esforços dos desenvolvedores~\cite{Bittar, Kleppe:2003}


No desenvolvimento tradicional, ou seja, sem seguir o MDD, artefatos de alto nível (ex., modelos, diagramas) usualmente são produzidos antes da codificação e costumam ser úteis apenas nas etapas iniciais do ciclo de desenvolvimento. Conforme o desenvolvimento evolui, mudanças são aplicadas somente no código-fonte e não nos modelos/diagramas. Assim, tais artefatos acabam se tornando incoerentes, ou seja, não refletem o que o código-fonte apresenta, o que faz com que o tempo e os esforços gastos na construção desses artefatos não sejam diretamente aproveitados na produção do software~\cite{Bittar}. Contrapartida, o MDD tem o foco nos modelos e busca simplificar o processo de desenvolvimento de software. Modelos são mais intuitivos para representação do conhecimento e menos dependentes do código-fonte, de forma que podem ser reutilizados facilmente em diferentes projetos. Diferentemente, o código-fonte possui uma linguagem que é densa e codificada, tornando-se difícil identificar, extrair e reutilizar o conhecimento apenas pela leitura do mesmo~\citep{Kleppe:2003}.



% subsection desenvolvimento_dirigido_a_modelos (end)

	 \subsection{Estudo sobre Modernização Dirigida à Arquitetura} % (fold)
	 \label{sub:estudo_sobre_moderniza_o_dirigida_arquitetura}
	 	
	 	

Reengenharia de Software e MDD são duas abordagens diferentes e que foram pesquisadas separadamente durante anos. Porém, recentemente pesquisadores identificam interesses análogos entre ambas abordagens. 

Como dito anteriormente, o principal objetivo da Reengenharia de Software é converter sistemas legados para novos sistemas sem alterar a sua funcionalidade, ou seja, isso implica em entender como o sistema legado foi implementado. Portanto, tal sistema legado deve ser convertido a um nível maior de abstração com o intuito de auxiliar o engenheiro de software a se concentrar apenas em informações importantes desse sistema, nesse ponto que MDD e Reengenharia de Software podem ser utilizados em conjunto.

Nesse contexto, a OMG OMG (\textit{Object Management Group}) propôs a abordagem Modernização Dirigida à Arquitetura (\textit{Architecture-Driven Modernization} - ADM), a qual tem como intuito automatizar e formalizar/padronizar o problema tradicional da Reengenharia de Software, i.e., ADM melhora a abordagem  de reengenharia de software tradicional com a utilização de MDD~\cite{PerezCastillo:2011jo}. Segundo~\citet{rezCastillo:2011gm} ADM é uma padronização definida pela OMG para auxiliar a atividade de reengenharia de sistemas legados, com o diferencial da utilização dos princípios da MDD (ver Seção~\ref{Cap2_Sec2_Desenvolvimento_Dirigido_a_Modelos}). ADM difere das abordagens tradicionais de reengenharia de software por dois principais motivos: (\textit{i}) ADM considera todos os artefatos de um sistema legado como modelos e (\textit{ii}) refatorações do sistema legado são realizadas nos modelos e depois gera-se um novo código-fonte refatorado tendo como base tais modelos. 

De acordo com~\citet{PerezCastillo:2011jo} ADM tem algumas vantagens quando comparada com abordagens de reengenharia convencionais: (\textit{i}) permite que pesquisadores definam técnicas de refatorações independente de linguagem e plataforma; (\textit{ii}) refatorações podem ser definidas como modelos, assim, podem ser reutilizadas, ou seja, utiliza o metamodelo KDM, o qual tem como um dos objetivos aumentar a interoperabilidade entre as ferramentas de reengenharia e (\textit{iii}) com a abordagem ADM é possível criar técnicas de refatorações genéricas e específicas. 

% subsection estudo_detalhado_sobre_moderniza_o_dirigida_arquitetura_ (end)

\begin{itemize}

\item estudar ferramentas para criar uma Linguagem Específica de Domínio: fez-se necessário identificar e estudar algumas ferramentas que auxiliam o desenvolvimento de Linguagem Específica de Domínio. Foram identificadas as seguintes ferramentas, XText\footnote{http://www.eclipse.org/Xtext/}, Eclipse EMF\footnote{http://www.eclipse.org/modeling/emf/} e Eclipse GMF\footnote{http://www.eclipse.org/modeling/gmp/}. Tais ferramentas serão estudadas para o desenvolvimento de um Linguagem Específica de Domínio, a qual será utilizada no terceiro passo da abordagem proposta;

\item Estudo detalhado da abordagem denominada Arquitetura Dirigida a Modelo - (ADM): conforme descrito no projeto submetido anteriormente pretende-se criar uma abordagem para auxiliar o engenheiro de software durante a atividade de reestruturação de sistemas legados. Dessa forma, vez se necessário estudar realizar a modernização de sistemas legados com a utilizanção de 

\end{itemize}

	 
	 	

	
	\subsection{Organização da Proposta}\label{org}

\section{Fundamentação Teórica}\label{sec:fundamentacao}


\subsection{Mineração de Padrões de Projeto}\label{sub:mineracao}


 
\subsection{Mineração de Interesses Transversais}\label{sub:mineracao}


\section{Proposta}\label{sec:proposta}


\subsection{Objetivos}\label{sec:Objetivos}


\subsection{Desafios de Pesquisa com Relação ao Projeto}\label{desafios}

\subsection{Metodologia}\label{sec:metodologia}

   
\subsubsection{Atividades e Cronograma}\label{sec:atividadeECronograma}


\subsection{Avaliação}\label{sec:avaliacao}

\subsection{Resultados Esperados}\label{sec:ResultadosEsperados}

\section{Trabalhos Relacionados}\label{sec:relacionados}
 
     

\singlespacing 
\bibliographystyle{icmc}
\bibliography{referencias2}

\end{document}
